\uuid{D4eZ}
\titre{Raisonnements}
\auteur{bodin ; croizat ; sacre}
\organisation{exo7}
\licence{CC BY-NC-SA 4.0}
\datecreate{2022-03-08}
\niveau{L1}
\module{Algèbre}
\chapitre{Logique, ensemble, raisonnement}
\sousChapitre{Absurde et contraposée}
\difficulte{2}
\theme{Raisonnements}
\exo7qcmid{hilisit.logique.045}

\contenu{
\begin{question}
\qclassification{100.03}
On souhaite prouver par récurrence, pour tout $n\ge0$, une proposition $\mathcal{P}_n$.
Après avoir prouvé $\mathcal{P}_0$, quelle rédaction du démarrage de l'étape d'hérédité convient ?
\begin{answers}
	\bad{Soit $n\ge0$. Je prouve $\mathcal{P}_1$.}
	\good{Soit $n\ge0$. Je suppose $\mathcal{P}_n$ vraie et je montre $\mathcal{P}_{n+1}$.}
	\bad{Soit $n\ge0$. Je suppose $\mathcal{P}_n$ vraie pour tout $n$ et je montre $\mathcal{P}_{n+1}$.}
	\bad{Soit $n\ge0$. Je suppose $\mathcal{P}_{n+1}$ vraie et je montre $\mathcal{P}_{n}$.}
\end{answers}
\begin{explanations}
La rédaction d'une récurrence est assez figée. L'étape d'hérédité commence toujours par "Soit $n\ge0$. Je suppose $\mathcal{P}_n$ vraie et je montre $\mathcal{P}_{n+1}$."
\end{explanations}
\end{question}
}
